\section{Exact computation and the economic accuracy contract}\label{sec:study61}
The preceding experiments establish distinct parts of the NBO procedure. We now examine how the exact implementation behaves away from one-survivor screening, what it costs inside a complete economic service, and whether the policies actually returned can be certified against the original Bellman optimum. The original nonlinear investment family and its continuous laws are retained. The additional two-control calculation is identified explicitly as a nested extension rather than being counted as another scalar-action policy service.

\subsection{Design, information sets, and accounting}
The completed comparison has four generators: a common-action controller with no fitted continuation, a quadratic continuation with native exact witnesses, a trained ReLU continuation with lossless screening, and the same trained ReLU construction with native finite-difference witnesses. The complete services use $(d,T)=(2,2),(4,4),(8,6)$, training seeds 60103 and 60209, and targets $q=9/10$ and $q=4/5$. The target requires expected policy cost at most $q$ times the expected cost of installed zero investment under the declared original initial law. These targets are not a claim about the optimum at every state.

Every service starts from the primitives. Its declared construction ladder uses $(256,32,4)$ and $(1024,64,8)$ for training samples, acquired leaves and innovation bins. The original action quantum is $1/4096$. Verification uses the same signed reference-advantage calculation, complete-cell feasibility rule and continuous-action lower cover for all four generators. Sequential inference examines 4096, 16,384 and 65,536 original-law interval paths, stopping at the first declared condition. The current execution has two hosted-worker repetitions of every mathematical design. Repeating identical seeds, policies and paths on a second worker does not double the number of independent economic observations.

The source freeze precedes the scientific services. A first execution failed in the parent recorder after an exclusive-write file was opened a second time. Its completed child outputs and failure logs remain in the attempt archive. A separately frozen correction changes that recorder to immutable numbered checkpoints and a single final ledger; the mathematical child service, seeds, targets and stopping rules are unchanged. The full corrected catalogue is used below, with no mixture of partial first-run clocks and completed rerun clocks. The separate centered-Bellman amendment was written after the uncentered frontier was observed. The fixed-policy revalidation and native factorial protocols likewise identify which predecessor results were known before their own frozen computations.

\subsection{Complete services and actual policy costs}
Table~\ref{tab:services61} reports every task--target--generator group on the complete return boundary. It includes imports, fitting, the executed construction and search attempts, whole-cell verification, sequential inference, durable scientific outputs, the process log and the first durable parent receipt. Cold use additionally charges native compilation; warm use starts with the compiled binary available. The source registry gives every one of the \NBOServices{} services, including the \NBOExhausted{} exhausted budgets. In total, \NBOAttained{} services meet their declared economic target.

\begin{table}[htbp]
\centering
\caption{Complete return work at the same economic targets}\label{tab:services61}
{\small\setlength{\tabcolsep}{3pt}
\begin{tabular}{lcccrrrr}
\hline
$d/T$ & $q$ & Mode & Met 0/1 & Warm 0 (s) & Warm 1 (s) & Cold 0 (s) & Cold 1 (s) \\
\hline
2/2 & 9/10 & C & 2/2 & 1.244 & 1.247 & 1.244 & 1.247 \\
2/2 & 9/10 & Q & 2/2 & 1.310 & 1.310 & 8.615 & 6.046 \\
2/2 & 9/10 & S & 2/2 & 1.562 & 1.573 & 1.562 & 1.573 \\
2/2 & 9/10 & N & 2/2 & 1.532 & 1.537 & 8.837 & 6.273 \\
2/2 & 4/5 & C & 2/2 & 3.792 & 3.793 & 3.792 & 3.793 \\
2/2 & 4/5 & Q & 2/2 & 3.981 & 3.956 & 11.286 & 8.692 \\
2/2 & 4/5 & S & 2/2 & 4.820 & 4.818 & 4.820 & 4.818 \\
2/2 & 4/5 & N & 2/2 & 4.781 & 4.730 & 12.086 & 9.466 \\
4/4 & 9/10 & C & 2/2 & 4.347 & 4.366 & 4.347 & 4.366 \\
4/4 & 9/10 & Q & 2/2 & 4.738 & 4.712 & 12.044 & 9.448 \\
4/4 & 9/10 & S & 2/2 & 5.523 & 5.558 & 5.523 & 5.558 \\
4/4 & 9/10 & N & 2/2 & 5.455 & 5.464 & 12.760 & 10.200 \\
4/4 & 4/5 & C & 2/2 & 15.835 & 15.887 & 15.835 & 15.887 \\
4/4 & 4/5 & Q & 2/2 & 17.012 & 16.904 & 24.317 & 21.640 \\
4/4 & 4/5 & S & 2/2 & 19.796 & 19.709 & 19.796 & 19.709 \\
4/4 & 4/5 & N & 2/2 & 19.524 & 19.858 & 26.830 & 24.594 \\
8/6 & 9/10 & C & 2/2 & 16.858 & 16.787 & 16.858 & 16.787 \\
8/6 & 9/10 & Q & 2/2 & 18.371 & 18.271 & 25.676 & 23.007 \\
8/6 & 9/10 & S & 2/2 & 19.675 & 19.631 & 19.675 & 19.631 \\
8/6 & 9/10 & N & 2/2 & 19.575 & 19.365 & 26.880 & 24.101 \\
8/6 & 4/5 & C & 0/0 & 52.701 & 52.132 & 52.701 & 52.132 \\
8/6 & 4/5 & Q & 0/0 & 57.282 & 57.480 & 64.588 & 62.216 \\
8/6 & 4/5 & S & 0/0 & 62.443 & 62.636 & 62.443 & 62.636 \\
8/6 & 4/5 & N & 0/0 & 61.769 & 61.469 & 69.075 & 66.206 \\
\hline\end{tabular}}
\par\smallskip
{\footnotesize C uses the common verified action family without a fitted continuation; Q uses native quadratic witnesses; S and N use the screened and native ReLU witnesses. The target is expected policy cost at most q times zero-policy cost. Each worker has two declared seeds; Met 0/1 gives the attained counts separately, each out of two. Warm and cold entries are within-worker two-seed medians of the complete durable-return clock. Cold work additionally charges that worker's measured native compilation for Q and N. Worker repetitions are not new independent economic samples. Failed targets remain in the table.}
\end{table}

\begin{table}[htbp]
\centering
\caption{Complete-service effect of replacing screened with native ReLU search}\label{tab:matched61}
{\small\setlength{\tabcolsep}{3pt}
\begin{tabular}{lcccl}
\hline
$d/T$ & $q$ & Saving 0 (\%) & Saving 1 (\%) & Faster 0/1 \\
\hline
2/2 & 9/10 & 1.95 & 2.29 & 2/2 \\
2/2 & 4/5 & 0.81 & 1.83 & 2/2 \\
4/4 & 9/10 & 1.24 & 1.68 & 2/2 \\
4/4 & 4/5 & 1.37 & -0.75 & 2/1 \\
8/6 & 9/10 & 0.51 & 1.35 & 1/2 \\
8/6 & 4/5 & 1.08 & 1.86 & 2/2 \\
\hline\end{tabular}}
\par\smallskip
{\footnotesize Savings are paired relative reductions in complete warm return time, summarized by the median of the two seeds within each worker. Negative values mean slower native service. Faster counts are out of two per worker. The exact returned policy, fitted parameters, stopping stage, inference looks and intervals coincide. These are descriptive timings, not a frequency-controlled or cross-platform speed theorem.}
\end{table}


The screened and native ReLU configurations return the same fitted parameters, action witnesses, acquired policy, verification arrays, stopping stage, path counts and cost intervals. Their difference is therefore an implementation-work difference at the same mathematical answer. The full-service comparison charges the native request and response as well as the exact computation. The native query benchmark below is not substituted for this complete return boundary, and a reduction in exact comparison counts is not presumed to translate proportionally into service time.

A different generator requires a different comparison. Table~\ref{tab:cost61} reports native-ReLU-minus-common and native-ReLU-minus-quadratic cost intervals for each task, target and seed on a fresh final-policy comparison stream. Positive intervals identify a more costly native-ReLU policy; exact identities are reported as identities; other intervals containing zero remain unresolved. The matched native-minus-screened contrast is exactly zero because their complete actor identities coincide. Conventional methods receive the same verifier and economic criterion. Their lower construction or return work is not removed from the ranking merely because they do not use learned ReLU features.

\begin{table}[htbp]
\centering
\caption{Actual returned-policy cost: native ReLU minus conventional comparators}\label{tab:cost61}
{\small\setlength{\tabcolsep}{3pt}
\begin{tabular}{lcrcc}
\hline
$d/T$ & $q$ & Seed & N minus C & N minus Q \\
\hline
2/2 & 9/10 & 60103 & [0.00000, 0.00000] & [0.00000, 0.00000] \\
2/2 & 9/10 & 60209 & [-0.00514, 0.00515] & [-0.00514, 0.00515] \\
2/2 & 4/5 & 60103 & [-0.00517, 0.00515] & [-0.00517, 0.00515] \\
2/2 & 4/5 & 60209 & [-0.00545, 0.00488] & [-0.00545, 0.00488] \\
4/4 & 9/10 & 60103 & [-0.00781, 0.00751] & [-0.00781, 0.00751] \\
4/4 & 9/10 & 60209 & [0.00000, 0.00000] & [0.00000, 0.00000] \\
4/4 & 4/5 & 60103 & [-0.00769, 0.00760] & [-0.00769, 0.00760] \\
4/4 & 4/5 & 60209 & [0.00000, 0.00000] & [0.00000, 0.00000] \\
8/6 & 9/10 & 60103 & [0.00000, 0.00000] & [0.00000, 0.00000] \\
8/6 & 9/10 & 60209 & [0.00000, 0.00000] & [0.00000, 0.00000] \\
8/6 & 4/5 & 60103 & [0.00000, 0.00000] & [0.00000, 0.00000] \\
8/6 & 4/5 & 60209 & [0.00000, 0.00000] & [0.00000, 0.00000] \\
\hline\end{tabular}}
\par\smallskip
{\footnotesize Intervals use the fresh final-policy comparison stream, conditional on the original selection. Endpoints are rounded outward. Positive intervals identify higher native-ReLU cost; an interval containing zero is unresolved. Exact policy identities have an exact zero interval. There are twelve mathematical designs, each repeated on two workers with the same policies and stream, not twenty-four independent comparisons. N minus S is identically zero.}
\end{table}


Across the twelve mathematical designs, native ReLU and each of the common-action and quadratic comparators have seven exact policy identities and five unresolved cost differences. No remaining comparison identifies a strictly lower native-ReLU cost. The common-action service has the smallest within-worker median warm return time in every task--target group; in a group whose target is exhausted by all generators this is a statement about observed work, not the work required to attain the target. These outcomes distinguish the implemented NBO contract from a claim of universal neural superiority.

The additional exact witness is available only after the common action family and the finite fitted menu have been evaluated. Comparing the selected policy immediately before and after that addition identifies the decisions attributable to the extra witness under a fixed verifier. The complete reconstruction checks every such cell, including unchanged and rejected proposals. Across the distinct reconstructed verification objects, it records \NBOExtraChanges{} additional-witness cell changes. This count includes the tested fitted generators; it is not an independent sample size, and it does not imply an identified average welfare gain for each changed cell. Actual economic consequences are assessed by the full-policy intervals rather than by assigning a policy value to a local reduction in a certificate.

\subsection{A language- and representation-matched action-search experiment}
The eight-method stress experiment varies three training-style fixture seeds, two feature widths, three action-lattice cardinalities and five regimes. The regimes include mixed activation, activation crossing, large signed cancellation, an exact tie between neighboring actions and a flat objective. Every rational objective is checked against exhaustive Fraction evaluation with the original smallest-index tie rule. An additional fixed subset receives seven shuffled repetitions per regime and worker. Coarser endpoint bounds and large-lattice guard cases are retained whether or not pruning is effective.

The eight methods separate symbolic reduction, vectorized enclosure, screening, exhaustive exact evaluation, native finite differences and adaptive subdivision. However, one direct contrast in that catalogue combines two changes: native exhaustive search uses the unreduced objective, whereas the native piecewise solver first reduces activation regions. We therefore execute a further two-by-two comparison before drawing an algorithm-specific timing conclusion. All four cells use the identical C++ arbitrary-size rational engine, objective, lattice, exact tie rule and original-value postcheck. They cross activation reduction, off or on, with exhaustive or finite-difference search. Reduction is charged to both reduced implementations. The same ninety fixed rational instances, rather than an outcome-selected favorable subset, are used on two newly allocated worker jobs.

\begin{table}[htbp]
\centering
\caption{Finite-difference versus exhaustive search at the same reduced representation}\label{tab:native-factorial61}
{\small\setlength{\tabcolsep}{3pt}
\begin{tabular}{rrrrrr}
\hline
Width & Lattice & Faster 0 & Faster 1 & Saving 0 (\%) & Saving 1 (\%) \\
\hline
8 & 33 & 13/15 & 8/15 & 2.51 & 0.40 \\
8 & 129 & 15/15 & 15/15 & 23.05 & 25.20 \\
8 & 513 & 15/15 & 15/15 & 60.51 & 65.07 \\
32 & 33 & 1/15 & 0/15 & -17.79 & -18.93 \\
32 & 129 & 15/15 & 15/15 & 26.77 & 30.69 \\
32 & 513 & 15/15 & 15/15 & 70.43 & 72.56 \\
\hline\end{tabular}}
\par\smallskip
{\footnotesize Both methods use the same C++ arbitrary-size rational engine, identical reduced objective and original-value postcheck. Each row includes three fixed seeds and all five regimes. Savings are medians of paired per-instance relative timings, not ratios of unrelated medians. Reduction, process creation, serialization and parsing are charged to both methods. The unreduced pair and all four-cell repeated-query results are retained in the supplement.}
\end{table}


At a 513-point lattice, finite differences save a median of approximately 61--73 percent relative to native exhaustive search at the same reduced representation, depending on width and worker. Every instance in those cells is faster. At width 32 and lattice size 33, the median saving is negative on both workers, approximately $-18$ and $-19$ percent. Thus the operation-count advantage at large lattices has an observed small-lattice boundary; a root-free exact solver is not uniformly the cheapest implementation. The separate width-eight, 33-point cell is close to the process-overhead floor and is more sensitive to the worker and instance. The complete four-cell repetition table retains these effects rather than attributing a language change or activation simplification to screening alone.

Batching supplies another measured distinction. The complete ninety-query catalogue is evaluated at batch sizes one, eight and thirty-two for each native implementation. Preprocessing, reduction where applicable, serialization, process creation, parsing and exact original-objective checks are all included. Every batch returns the same ninety exact answers. The batch study changes only the boundary of an implementation service; it does not create a new fitted policy or a new economic stopping experiment. Its full timings and byte counts are reported in the supplement, separately from the complete policy-service clocks.

Flatness and cancellation also clarify what a lossless result means. An interval screen is allowed to retain every minimizer and many nonminimizers; it need not produce one survivor. The finite-difference construction instead uses the exact polynomial piece, retaining the first minimizer even on a flat segment. Adaptive subdivision can reach its node budget and call an exact fallback. A fallback preserves correctness but charges its prefix work in addition to the subsequent solver. The full stress tables report all fallback events, endpoint resolutions, recorded operand lengths and interval-array memory. Endpoint coarsening after an exact rational bound is not described as a change in internal floating-point precision.

\subsection{Two controls and an additional innovation}
The multi-action calculation uses saved, genuinely trained two-state ReLU continuations at two specified states. The two controls share capacity $a_1+a_2\le1/8$. Their exposure columns are $(1/2,1/4)$ and $(1/4,1/2)$, and the optional second independent shock has a different exposure. Quantum denominators 64 and 128 give 45 and 153 feasible lattice points. Setting the second action to zero and disabling the second shock recovers the original action-dependent objective.

\begin{table}[htbp]
\centering
\caption{Constrained two-control exact witness recovery}\label{tab:multi61}
{\small\setlength{\tabcolsep}{3pt}
\begin{tabular}{rrrrrrrr}
\hline
Cap index & Shocks & Designs & Exhaustive & Adaptive & Max nodes & Fallback & Time ratio \\
\hline
8 & 1 & 4 & 45 & 7.0 & 107 & 0 & 3.24 \\
8 & 2 & 4 & 45 & 7.0 & 107 & 0 & 3.19 \\
16 & 1 & 4 & 153 & 80.0 & 256 & 4 & 2.70 \\
16 & 2 & 4 & 153 & 80.0 & 256 & 4 & 2.68 \\
\hline\end{tabular}}
\par\smallskip
{\footnotesize Each of four mathematical designs per row uses a saved genuinely fitted ReLU continuation, two state points and two training seeds, repeated on two workers. Shared capacity is 1/8, with quantum denominator eight times Cap index. Adaptive counts are selected exact evaluations; node-bound computation is additional work. Time ratio is the median paired adaptive/exhaustive query time over the eight records. A ratio above one means slower adaptive execution despite fewer exact evaluations. These are action-search results, not complete two-control welfare comparisons.}
\end{table}


Both exact exhaustive evaluation and verified subdivision return the original constrained minimizer in every recorded query. The analytic uniform-box formula includes all shock dimensions actually used, with signed fitted output weights. Fewer exact point evaluations do not eliminate the cost of computing rectangle bounds: the adaptive/exhaustive time ratio is therefore reported along with evaluations and visited nodes. These experiments validate and cost a non-scalar witness extension. They do not constitute a complete two-control policy-cost or welfare study, and the scalar quartic search theorem is not invoked for a two-shock objective outside its stated degree structure.

\subsection{Attaining an original Bellman-accuracy target}
The direct Bellman calculation retains the original two-state, two-date economy. Lower intervals cover every continuous feasible action, whereas upper actions are quantum-valued and feasible throughout each acquired cell. Terminal innovation integration is analytic; unresolved earlier expectations use enclosures over the full original innovation bins. The uncentered procedure gives a valid but wide bracket. Centering at the same zero-policy continuation on both sides preserves cancellation before the gap is formed.

\begin{table}[htbp]
\centering
\caption{Original-model whole-domain Bellman accuracy}\label{tab:bellman61}
{\small\setlength{\tabcolsep}{3pt}
\begin{tabular}{rrrrrr}
\hline
$N$ & $A$ & $q$ & Uncentered gap & Centered gap & Centered work (s) \\
\hline
8 & 8 & 4 & 2.18231 & 0.62981 & 0.478 \\
16 & 16 & 8 & 1.12578 & 0.28126 & 2.290 \\
32 & 32 & 16 & 0.59109 & 0.14409 & 11.288 \\
64 & 64 & 32 & 0.30286 & 0.07289 & 80.081 \\
128 & 64 & 32 & -- & 0.03979 & 289.742 \\
\hline\end{tabular}}
\par\smallskip
{\footnotesize The model remains the original two-state, two-date investment economy with its continuous feasible actions and innovation law. N is the number of verification cells per state coordinate. The common-reference lower and upper endpoints cancel the same reference value pointwise. Centered work includes every preceding declared rung and endpoint serialization on that execution. The uncentered gap is the corresponding separately executed numerical bracket, not a same-host timing comparator. Tolerances 1/2, 1/4, 1/8 and 1/16 are attained; 1/32 is not.}
\end{table}


The centered ladder reaches the all-state tolerances $1/2$, $1/4$, $1/8$ and $1/16$. At $N=128$, its maximum datewise gap is at most \NBOCenteredGap{}, compared with approximately $0.30330$ at the final uncentered rung. The centered construction does not attain $1/32$ on its declared ladder. Its cumulative recorded work includes every earlier centered rung and serialized endpoint array. These are constructive bounds against the original continuous-action optimum, not confidence intervals for gains relative to zero investment. The state partition here is a tensor verification grid; the non-tensor prospective actor and its high-dimensional service evidence remain distinct objects.

The bracket by itself certifies the policy constructed by that Bellman calculation. To assess the policies actually returned by the prospective services, we separately execute Proposition~\ref{prop:fixed-policy61}. The complete two-date service catalogue has \NBOFixedAliases{} returned records and \NBOFixedActors{} distinct policy-plus-partition identities. Every identity is included, whether its policy came from a neural, quadratic or common-action generator. The lower table is fixed in advance at $N=128$, $A=64$, $q=32$. The upper recursion is recomputed for each original actor through all intersecting acquired leaves. Neither training nor action selection is repeated.

\begin{table}[htbp]
\centering
\caption{Revalidation of every distinct actually returned two-date actor}\label{tab:fixed-policy61}
{\small\setlength{\tabcolsep}{3pt}
\begin{tabular}{lrcrrrrr}
\hline
Actor & Leaves & Modes & Aliases & Date 0 gap & All-date gap & Unclipped & Own work (s) \\
\hline
\texttt{2cc41f05} & 64 & N/S & 4 & 0.07176 & 0.07176 & 0.07176 & 2.254 \\
\texttt{3f2668fe} & 64 & C/Q & 8 & 0.07176 & 0.07176 & 0.07176 & 2.278 \\
\texttt{4e5b6e26} & 32 & N/S & 4 & 0.23801 & 0.30330 & 0.30330 & 2.228 \\
\texttt{718597a7} & 64 & N/S & 4 & 0.06778 & 0.06778 & 0.06778 & 2.279 \\
\texttt{a06e05ba} & 32 & C/N/Q/S & 12 & 0.23801 & 0.30330 & 0.30330 & 2.223 \\
\hline\end{tabular}}
\par\smallskip
{\footnotesize All thirty-two returned service records are represented by five complete policy-plus-partition identities. No actor is retrained, improved or replaced. Verification uses N=128 and q=32 with the deposited centered lower Bellman table. Closed actor ranges include every intersecting leaf. The independent global nonworsening bound is intersected with upper excess endpoints, but the maximum gaps here coincide with the retained unclipped calculation. Own work covers each fixed-actor recursion through its endpoint archive; the shared lower-table reconstruction and process-return overhead are separate audit records.}
\end{table}


Three of these returned actors have all-state bounds between approximately $0.06778$ and $0.07176$, and attain tolerance $1/8$ but not $1/16$. The other two retain bounds around $0.30330$. The policy distinction is economically important: a service that meets a ten-percent initial-law improvement target can still return an actor with a materially wider worst-state optimality account. Conversely, the tighter fixed-policy bounds establish original-optimum accuracy for actual deployed actors rather than assigning the Bellman comparator's bound to an unrelated neural policy. The maximum gap is unchanged by the independently justified zero intersection in these results; the unclipped endpoints and gaps are also retained.

The shared lower-table reconstruction, each actor's additional upper recursion and the whole-process receipt are charged separately. These post-selection verification costs are not spliced into a fictitious earlier stopping run, and no old service clock is replaced by a replay timing. Deterministic complete-domain verification adds no independent policy-cost observation.

\subsection{Economic interpretation and computational scope}
For two implementations returning the same policy, the gain is a reduction in computation cost, conditional on the measured execution environment. The setup-and-use formula above gives the full range of implementation charges and use counts supported at any stated price of the work unit. For different policies, a returned-policy gain interval $[l,u]$, a replacement fee $\tau$ and an incremental measured-work price $\lambda\Delta w$ give the net-benefit interval
\[
 [l-\tau-\lambda\Delta w,\ u-\tau-\lambda\Delta w].
\]
A positive lower endpoint supports replacement under that price specification; a negative upper endpoint rules it out; the remaining cases are unresolved. This algebra does not estimate an installation fee, an electricity price or a welfare normalization from a hosted-runner clock.

The experiments specify CPU affinity and record the available processor, compiler and runtime metadata. Processor frequency is not controlled and hardware counters were unavailable under the recorded permissions. Separate worker allocations do not establish independent physical architectures. Rational operation counts, observed operand-size diagnostics, interval-array bytes, process resident memory, cold compilation and complete return times therefore remain separate quantities. The new evidence supports a terminating, costed exact witness and an executable original-optimum certificate within the original NBO development. It does not require suppressing adverse conventional comparisons or inferring an unexecuted scaling result.
